%% toc: P3 39 39.6
\chapter{Least and Partial Fixed-Point Logics}
\label{sec:p03-w03-p03-059-fixed-point-logics}

First-order logic cannot define the transitive closure of a relation on
finite structures. Fixed-point logics repair this by letting a formula
define a relation through iteration: a formula \(\varphi(R,\bar x)\) with a
free relation variable \(R\) is read as an operator that sends a relation to
a new relation, and the logic adds a constructor naming a fixed point of
that operator. This chapter develops the two fixed-point logics that matter
most for descriptive complexity. Least fixed-point logic (LFP) iterates only
operators made monotone by a positivity condition. Partial fixed-point logic
(PFP) iterates arbitrary first-order operators on finite structures and
returns the empty relation when the iteration never stabilizes.

The inputs are a finite relational vocabulary \(\tau\), first-order
formulas over \(\tau\) together with one extra relation variable, and, for
the complexity results, finite structures equipped with a linear order. The
conclusions are: the operator theory that makes least and greatest fixed
points exist (Knaster--Tarski and the stage construction); the syntax and
semantics of both logics; Immerman's normal form, by which on finite
structures negated least fixed points can be eliminated, proved through
stage comparison; the inclusion of LFP in PFP; the fact that a PFP
iteration can run through exponentially many stages; the evaluation of PFP
formulas in polynomial space; the capture of \(\mathsf{PSPACE}\) by PFP on
ordered finite structures; and the resulting equivalence between the
question \(\mathsf P=\mathsf{PSPACE}\) and the question whether LFP and PFP
define the same classes of ordered finite structures.

The presentation follows Gr\"adel's chapter in the author-hosted PDF of
2007, whose running page numbers are used for every locator below
\cite[pp.~152--180]{Gradel2007FMTDescriptiveChapter}; whether that file coincides with the
printed Springer text of the same year has not been checked. The handbook's
own first-order least fixed-point constructor
(\cref{def:p04-fol-lfp}) is the same object as the constructor of
\cref{def:p03-w03-p03-059-lfp}; the present chapter adds the operator
theory behind it, the greatest fixed points, and the partial fixed points.
Complexity classes are those of \cref{def:p04-resource-classes}, and
``captures'' has the meaning of \cref{def:p04-captures-complexity}.

\section{Operators and fixed points on power-set lattices}
\label{sec:p03-w03-p03-059-operators}

Gr\"adel states the basic fixed-point theory for the power-set lattice
\((\mathcal P(B),\subseteq)\) of an arbitrary set \(B\), finite or infinite
\cite[p.~153]{Gradel2007FMTDescriptiveChapter}. In applications \(B\) is \(A^k\), the set
of \(k\)-tuples over the universe \(A\) of a structure.

\begin{hbdefinition}{def:p03-w03-p03-059-operator}{Fixed points and monotone operators}
Let \(B\) be a set and \(F:\mathcal P(B)\to\mathcal P(B)\) a function,
called an \emph{operator} on \(B\).
\begin{enumerate}
 \item A set \(X\subseteq B\) is a \emph{fixed point} of \(F\) when
       \(F(X)=X\).
 \item A \emph{least fixed point} of \(F\) is a fixed point \(X\) with
       \(X\subseteq Y\) for every fixed point \(Y\) of \(F\); a
       \emph{greatest fixed point} is a fixed point \(X\) with
       \(Y\subseteq X\) for every fixed point \(Y\).
 \item \(F\) is \emph{monotone} when \(X\subseteq Y\) implies
       \(F(X)\subseteq F(Y)\) for all \(X,Y\subseteq B\).
\end{enumerate}
A least (greatest) fixed point is unique when it exists, since two of them
are each contained in the other. This is Definition~3.3.1 of
\cite[p.~153]{Gradel2007FMTDescriptiveChapter}.
\index{fixed point!of an operator}\index{monotone operator}
\end{hbdefinition}

An operator need not have any fixed point: on a nonempty \(B\), the
complement operator \(X\mapsto B\setminus X\) has none, because
\(X=B\setminus X\) is impossible. Monotonicity is the hypothesis that
removes this obstruction.

\begin{hbtheorem}{thm:p03-w03-p03-059-knaster-tarski}{proved=annex:p03-w03-p03-059-knaster-tarski}{Knaster--Tarski for power-set lattices}
Every monotone operator \(F:\mathcal P(B)\to\mathcal P(B)\) has a least
fixed point \(\operatorname{lfp}(F)\) and a greatest fixed point
\(\operatorname{gfp}(F)\), and
\[
\begin{aligned}
 \operatorname{lfp}(F)&=\bigcap\{X:F(X)=X\}=\bigcap\{X:F(X)\subseteq X\},\\
 \operatorname{gfp}(F)&=\bigcup\{X:F(X)=X\}=\bigcup\{X:X\subseteq F(X)\}.
\end{aligned}
\]
This is Theorem~3.3.2 of \cite[p.~153]{Gradel2007FMTDescriptiveChapter}.
\index{Knaster--Tarski theorem!power-set form}
\end{hbtheorem}

The intersection formula describes the least fixed point from above, as the
smallest set closed under \(F\). The stages below describe it from below,
by iterating \(F\) from the empty set.

\begin{hbdefinition}{def:p03-w03-p03-059-stages}{Stages, inductive operators and closure ordinal}
For an operator \(F\) on \(B\), the \emph{stages} \(X^\alpha\), for
ordinals \(\alpha\), are defined by
\[
 X^0\defeq\varnothing,\qquad X^{\alpha+1}\defeq F(X^\alpha),\qquad
 X^\lambda\defeq\bigcup_{\alpha<\lambda}X^\alpha\ \text{for limit }\lambda.
\]
\(F\) is \emph{inductive} when the stages increase:
\(X^\alpha\subseteq X^\beta\) whenever \(\alpha<\beta\). For an inductive
operator the stages reach a fixed point
\(X^\infty\) (\cref{thm:p03-w03-p03-059-closure-ordinal}), called its
\emph{inductive fixed point}, and the least ordinal \(\alpha\) with
\(X^\alpha=X^{\alpha+1}\) is the \emph{closure ordinal}
\(\operatorname{cl}(F)\). Dually, the \emph{downward stages} are
\(Y^0\defeq B\), \(Y^{\alpha+1}\defeq F(Y^\alpha)\) and
\(Y^\lambda\defeq\bigcap_{\alpha<\lambda}Y^\alpha\) at limits. For
\(X\subseteq B\) write \(\overline X\defeq B\setminus X\); the \emph{dual
operator} of \(F\) is \(F^d(X)\defeq\overline{F(\overline X)}\)
\cite[p.~154]{Gradel2007FMTDescriptiveChapter}.
\index{stages of an operator}\index{closure ordinal}\index{dual operator}
\end{hbdefinition}

\begin{hblemma}{thm:p03-w03-p03-059-closure-ordinal}{proved=annex:p03-w03-p03-059-closure-ordinal}{The stages of an inductive operator stop}
Let \(F\) be an inductive operator on \(B\). There is an ordinal
\(\alpha\) with \(X^\alpha=X^{\alpha+1}\); from that ordinal on the stages
are constant, and \(X^\infty=X^{\operatorname{cl}(F)}\) is a fixed point
of \(F\). Moreover \(|\operatorname{cl}(F)|\leq|B|\). Every monotone
operator is inductive. This is Lemma~3.3.3 of
\cite[p.~154]{Gradel2007FMTDescriptiveChapter} together with the observations before it.
\end{hblemma}

\begin{hbtheorem}{thm:p03-w03-p03-059-stages-lfp}{proved=annex:p03-w03-p03-059-stages-lfp}{Stage construction of the extremal fixed points}
For a monotone operator \(F\) on \(B\):
\begin{enumerate}
 \item the inductive fixed point is the least fixed point,
       \(X^\infty=\operatorname{lfp}(F)\);
 \item the downward stages decrease and become constant at a fixed point
       \(Y^\infty\), and \(Y^\infty=\operatorname{gfp}(F)\);
 \item \(F^d\) is monotone, and
       \(\operatorname{lfp}(F)=\overline{\operatorname{gfp}(F^d)}\) and
       \(\operatorname{gfp}(F)=\overline{\operatorname{lfp}(F^d)}\).
\end{enumerate}
Item~1 is Proposition~3.3.4; items~2 and~3 are stated on
\cite[p.~154]{Gradel2007FMTDescriptiveChapter}, item~3 being left there as
Exercise~3.3.5.
\end{hbtheorem}

\begin{hblemma}{thm:p03-w03-p03-059-finite-stages}{proved=annex:p03-w03-p03-059-finite-stages}{Finite stages reach the fixed points in polynomially many steps}
Let \(A\) be a finite set with \(n\) elements and
\(F:\mathcal P(A^k)\to\mathcal P(A^k)\) monotone. Then
\(\operatorname{lfp}(F)=X^m\) and \(\operatorname{gfp}(F)=Y^m\) for some
\(m\leq n^k\), so only finite stages are needed. If \(F\) is computable in
time polynomial in \(n\), then so are \(\operatorname{lfp}(F)\) and
\(\operatorname{gfp}(F)\). The complexity claim is Lemma~3.3.6 of
\cite[p.~155]{Gradel2007FMTDescriptiveChapter}.
\end{hblemma}

\section{Least fixed-point logic}
\label{sec:p03-w03-p03-059-lfp}

Relation variables, like first-order variables, have binding scopes. In the
polarity conditions and relation substitutions below, occurrences of a
relation symbol are those free with respect to relation binding; a nested
fixed-point binder for that symbol starts a new scope. Before substitution,
bound relation variables are renamed to distinct fresh symbols of the same
arity, avoiding the free symbols of the formulas being inserted. Bound
first-order variables are likewise renamed to prevent capture of parameters.
Such renaming changes neither the update operator nor the defined relation;
the annex records the semantic justification.

Let \(\tau\) be a relational vocabulary, \(R\) a relation symbol not in
\(\tau\) of arity \(k\), and \(\varphi(R,\bar x)\) a formula of vocabulary
\(\tau\cup\{R\}\) whose free first-order variables include the \(k\)
distinct variables \(\bar x\). In a \(\tau\)-structure \(\mathfrak A\) with
universe \(A\), and for a fixed assignment to the other free variables of
\(\varphi\), the formula defines the \emph{update operator}
\[
 F_\varphi:\mathcal P(A^k)\to\mathcal P(A^k),\qquad
 F_\varphi(R)\defeq\{\bar a:(\mathfrak A,R)\entails\varphi(R,\bar a)\}
\]
\cite[p.~152]{Gradel2007FMTDescriptiveChapter}. An occurrence of \(R\) is \emph{positive}
when it lies under an even number of negations, counted after implications
and biconditionals are rewritten with negation, conjunction and
disjunction. Positivity is the syntactic condition that makes
\(F_\varphi\) monotone, and so brings
\cref{thm:p03-w03-p03-059-knaster-tarski} into play.

\begin{hblemma}{thm:p03-w03-p03-059-positive-monotone}{proved=annex:p03-w03-p03-059-positive-monotone}{Positivity gives monotone operators}
If every occurrence of \(R\) in \(\varphi(R,\bar x)\) is positive, then for
every structure \(\mathfrak A\) and every assignment to the other free
variables, \(F_\varphi\) is monotone. If every occurrence is negative
(under an odd number of negations), \(F_\varphi\) is antitone:
\(S\subseteq S'\) implies \(F_\varphi(S')\subseteq F_\varphi(S)\).
\end{hblemma}

\begin{hbdefinition}{def:p03-w03-p03-059-lfp}{Least fixed-point logic}
LFP extends first-order logic by the following formation rule. If
\(\varphi(R,\bar x)\) is a formula of vocabulary \(\tau\cup\{R\}\) in which
\(R\) occurs only positively, \(\bar x\) is a \(k\)-tuple of distinct
variables, \(k\) is the arity of \(R\), and \(\bar t\) is a \(k\)-tuple of
terms, then
\[
 [\mathbf{lfp}\,R\bar x\,.\,\varphi](\bar t)
 \quad\text{and}\quad
 [\mathbf{gfp}\,R\bar x\,.\,\varphi](\bar t)
\]
are formulas of vocabulary \(\tau\). Their free first-order variables are
those of \(\varphi\) other than the variables of \(\bar x\), together with
the variables of \(\bar t\); \(R\) is bound. The free variables of
\(\varphi\) outside \(\bar x\) are the \emph{parameters} of the fixed-point
formula.

\emph{Semantics.} For a \(\tau\)-structure \(\mathfrak A\) and an
assignment to the free variables,
\(\mathfrak A\entails[\mathbf{lfp}\,R\bar x\,.\,\varphi](\bar t)\) exactly
when the tuple \(\bar t^{\mathfrak A}\) of values of \(\bar t\) belongs to
\(\operatorname{lfp}(F_\varphi)\), the operator being taken with the
parameters at their assigned values; the \(\mathbf{gfp}\) formula uses
\(\operatorname{gfp}(F_\varphi)\) instead. By
\cref{thm:p03-w03-p03-059-positive-monotone,thm:p03-w03-p03-059-knaster-tarski}
both fixed points exist on every structure, finite or infinite. This is
Definition~3.3.7 of \cite[p.~155]{Gradel2007FMTDescriptiveChapter}. The constructor of
\cref{def:p04-fol-lfp}, written there
\([\operatorname{lfp}_{R,\bar x}\varphi](\bar t)\), is the
\(\mathbf{lfp}\) case.
\index{least fixed-point logic}\index{LFP}\index{greatest fixed point!formula}
\end{hbdefinition}

\begin{hbexample}{def:p03-w03-p03-059-tc}{Transitive closure as a least fixed point}
For a binary relation symbol \(E\), Gr\"adel's Example~3.3.8 defines the
transitive closure by
\[
 [\mathbf{lfp}\,Txy\,.\,(Exy\lor\exists z(Exz\land Tzy))](u,v)
\]
\cite[p.~155]{Gradel2007FMTDescriptiveChapter}. The variable \(T\) occurs positively. On a
finite graph the stage \(T^m\) is the set of pairs joined by a nonempty
\(E\)-path with at most \(m\) edges: \(T^1=E\), and a pair enters
\(T^{m+1}\) exactly when it is an edge or its first edge is followed by a
path in \(T^m\). The union of the stages is the transitive closure, which
\cref{thm:p03-w03-p03-059-stages-lfp} identifies with the least fixed
point. The same example with \(u\) as a parameter is
\([\mathbf{lfp}\,Ty\,.\,(Euy\lor\exists x(Tx\land Exy))](v)\).
\end{hbexample}

\begin{hbtheorem}{thm:p03-w03-p03-059-gfp-dual}{proved=annex:p03-w03-p03-059-gfp-dual}{Greatest fixed points through negated least fixed points}
Let \(\varphi(R,\bar x)\) have only positive occurrences of \(R\), and let
\(\neg\varphi[R/\neg R]\) be the negation of the result of replacing every
free occurrence of \(R\bar s\) in \(\varphi\) by \(\neg R\bar s\). Then \(R\) occurs only
positively in \(\neg\varphi[R/\neg R]\), and on every structure
\[
 [\mathbf{gfp}\,R\bar x\,.\,\varphi](\bar t)
 \;\Leftrightarrow\;
 \neg[\mathbf{lfp}\,R\bar x\,.\,\neg\varphi[R/\neg R]](\bar t).
\]
This equivalence is displayed on \cite[p.~156]{Gradel2007FMTDescriptiveChapter}; it shows
that greatest fixed points add no expressive power once negation of
\(\mathbf{lfp}\) formulas is allowed.
\end{hbtheorem}

Because of \cref{thm:p03-w03-p03-059-gfp-dual}, unrestricted LFP formulas
may place negations in front of fixed points, and so may nest least and
greatest fixed points. The classical theory of inductive definitions works
with a narrower language.

\begin{hbdefinition}{def:p03-w03-p03-059-lfp1}{Positive least fixed-point logic}
\(\mathrm{LFP}_1\), also called positive LFP, is obtained from first-order
logic by forming least fixed points
\([\mathbf{lfp}\,R\bar x\,.\,\varphi](\bar x)\) of first-order formulas
\(\varphi\), positive in \(R\) and without parameters, and closing the
resulting formulas under conjunction, disjunction, and existential and
universal quantification, but not under negation
\cite[p.~173]{Gradel2007FMTDescriptiveChapter}. Gr\"adel refers to the preceding chapter
of the same book for a more formal definition; the description given here is
the one on p.~173.
\index{positive least fixed-point logic}\index{LFP1@\(\mathrm{LFP}_1\)}
\end{hbdefinition}

On finite structures these negations can be removed again. The proof below
follows Immerman's argument for the collapse of the fixed-point hierarchy
\cite[Thm.~2 and Lemmas~4.1--4.7, pp.~94--100]{Immerman1986}; it works with
an auxiliary class of formulas in which fixed points may be nested, may
carry parameters and may define several relations at once, and in which
negation stands only in front of atoms.

\begin{hbdefinition}{def:p03-w03-p03-059-simultaneous-lfp}{Simultaneous least fixed points and positive fixed-point formulas}
Let \(\sigma\) be a relational vocabulary; its symbols may include free
relation variables, which are then treated like any other relation symbol.
A \emph{system} \(\mathcal S\) over \(\sigma\) consists of distinct relation
symbols \(R_1,\ldots,R_m\notin\sigma\) of arities \(k_1,\ldots,k_m\)
and, for each \(i\), a \(k_i\)-tuple \(\bar x_i\) of distinct variables and
a formula \(\varphi_i\) of vocabulary \(\sigma\cup\{R_1,\ldots,R_m\}\),
written \(R_i\bar x_i\defeq\varphi_i\). The free variables of the
\(\varphi_i\) outside \(\bar x_i\) are the \emph{parameters} of
\(\mathcal S\). In a \(\sigma\)-structure \(\mathfrak A\), with the
parameters assigned, the system defines the operator
\[
 F_{\mathcal S}(X_1,\ldots,X_m)\defeq(Y_1,\ldots,Y_m),\qquad
 Y_i\defeq\{\bar a\in A^{k_i}:(\mathfrak A,X_1,\ldots,X_m)\entails\varphi_i(\bar a)\},
\]
on \(\mathcal P(A^{k_1})\times\cdots\times\mathcal P(A^{k_m})\). This
product is identified with the power set of the disjoint union
\(B\defeq\{(i,\bar a):i\leq m,\ \bar a\in A^{k_i}\}\), so that inclusion,
stages, \cref{thm:p03-w03-p03-059-knaster-tarski} and
\cref{thm:p03-w03-p03-059-finite-stages} apply componentwise.

The class \(\mathrm{Pos}[\sigma]\) of \emph{positive fixed-point formulas}
of vocabulary \(\sigma\) is the least class of formulas such that
\begin{enumerate}
 \item every \emph{literal} \(P\bar t\), \(\neg P\bar t\), \(s=t\),
       \(\neg s=t\) with \(P\in\sigma\) belongs to it;
 \item it is closed under \(\land\), \(\lor\), \(\exists x\) and
       \(\forall x\);
 \item if \(\mathcal S\) is a system over \(\sigma\) whose bodies
       \(\varphi_i\) belong to \(\mathrm{Pos}[\sigma\cup\{R_1,\ldots,R_m\}]\)
       and contain no literal \(\neg R_l\bar t\) for any \(l\leq m\), then
       \([\mathbf{lfp}\,R_j:\mathcal S](\bar t)\) belongs to it, for every
       \(j\leq m\) and every \(k_j\)-tuple \(\bar t\) of variables. Its free
       variables are those of \(\bar t\) and the parameters of
       \(\mathcal S\), and it holds when \(\bar t\) denotes a tuple in the
       \(j\)-th component of \(\operatorname{lfp}(F_{\mathcal S})\).
\end{enumerate}
The notation \([\mathbf{lfp}\,R_j:\mathcal S](\bar t)\) is that of
Gr\"adel's simultaneous fixed points \cite[Thm.~3.3.43,
p.~174]{Gradel2007FMTDescriptiveChapter}; the class \(\mathrm{Pos}[\sigma]\) is an
auxiliary class of this chapter. Write \(\bot\defeq\exists v\,\neg v=v\)
and \(\top\defeq\forall v\,v=v\); structures are nonempty, so these are
false and true. For \(\theta\in\mathrm{Pos}[\sigma\cup\{R\}]\), with \(R\)
of arity \(k\), and formulas \(\alpha(\bar u),\beta(\bar u)\) with a
\(k\)-tuple \(\bar u\) of free variables, \(\theta[R:\alpha/\beta]\) is the
result of replacing every free \(R\)-literal \(R\bar t\) by
\(\alpha(\bar t)\) and every free \(R\)-literal \(\neg R\bar t\) by
\(\beta(\bar t)\). First-order and relation variables bound in the formulas
are renamed so that no free variable of \(\alpha\) or \(\beta\) is captured;
occurrences under a binder for \(R\) are not substituted.
Two instances recur: \(\theta[R:\varnothing]\defeq\theta[R:\bot/\top]\),
and \(\theta^{\sim}\defeq\theta[R:\neg R\bar u/R\bar u]\), which exchanges
the two kinds of \(R\)-literals.
\index{simultaneous least fixed point}\index{positive fixed-point formula}
\end{hbdefinition}

The semantics in item~3 presupposes that \(F_{\mathcal S}\) is monotone, so
that \(\operatorname{lfp}(F_{\mathcal S})\) exists. The first lemma supplies
this, together with the substitution rule that the later constructions use
throughout.

\begin{hblemma}{thm:p03-w03-p03-059-pos-basic}{proved=annex:p03-w03-p03-059-pos-basic}{Monotonicity and substitution for positive fixed-point formulas}
\begin{enumerate}
 \item If a symbol \(P\in\sigma\) occurs in no literal \(\neg P\bar t\) of
       \(\theta\in\mathrm{Pos}[\sigma]\), then \(\theta\) is monotone in
       \(P\): if two \(\sigma\)-structures \(\mathfrak A,\mathfrak A'\)
       differ only in that \(P^{\mathfrak A}\subseteq P^{\mathfrak A'}\),
       then \(\mathfrak A\entails\theta[s]\) implies
       \(\mathfrak A'\entails\theta[s]\) for every assignment \(s\). In
       particular every system allowed in item~3 of
       \cref{def:p03-w03-p03-059-simultaneous-lfp} has a monotone operator,
       and the semantics there is well defined.
 \item Let \(\theta\in\mathrm{Pos}[\sigma\cup\{R\}]\), with \(R\notin\sigma\)
       of arity \(k\), and let \(\alpha,\beta\in\mathrm{Pos}[\sigma']\) with
       \(\sigma\subseteq\sigma'\), where the relation variables bound in
       \(\theta\) do not occur in \(\alpha,\beta\). Then
       \(\theta[R:\alpha/\beta]\in\mathrm{Pos}[\sigma']\). Let
       \(\mathfrak B\) be a \(\sigma'\)-structure with universe \(A\) and
       \(s\) an assignment, let \(Q\) be the set of \(\bar a\in A^k\) with
       \(\mathfrak B\entails\alpha(\bar a)\) under \(s\), and suppose that
       \(\beta\) defines \(A^k\setminus Q\) under \(s\). Then
       \(\mathfrak B\entails\theta[R:\alpha/\beta][s]\) if and only if
       \(\mathfrak B[R\mapsto Q]\entails\theta[s]\), where
       \(\mathfrak B[R\mapsto Q]\) interprets \(R\) as \(Q\) and agrees
       with \(\mathfrak B\) elsewhere. If \(\theta\) has no literal
       \(\neg R\bar t\), the hypothesis on \(\beta\) is not needed.
\end{enumerate}
\end{hblemma}

The decisive tool compares the stages at which tuples enter a least fixed
point.

\begin{hbdefinition}{def:p03-w03-p03-059-rank}{Rank and stage comparison}
Let \(F\) be a monotone operator on \(\mathcal P(A^k)\), \(A\) finite, with
stages \(X^0=\varnothing\) and \(X^{m+1}=F(X^m)\). The \emph{rank} of
\(\bar a\in A^k\) is \(|\bar a|_F\defeq\min\{m:\bar a\in X^m\}\) if
\(\bar a\in\operatorname{lfp}(F)\), and \(|\bar a|_F\defeq\infty\)
otherwise; every finite rank is at least \(1\). The \emph{stage comparison
relations} are
\[
 \bar a\le_F\bar b\iff|\bar a|_F\leq|\bar b|_F\text{ and }|\bar a|_F<\infty,
 \qquad
 \bar a<_F\bar b\iff|\bar a|_F<|\bar b|_F .
\]
These are Immerman's rank and his relations \(\leq_\varphi\) and
\(<_\varphi\) \cite[p.~97]{Immerman1986}, written here for the operator
rather than for the formula.
\index{rank of a tuple in a fixed point}\index{stage comparison}
\end{hbdefinition}

\begin{hblemma}{thm:p03-w03-p03-059-stage-comparison}{proved=annex:p03-w03-p03-059-stage-comparison}{Positive stage comparison and positive negation of least fixed points}
Let \(R\notin\sigma\) have arity \(k\), let \(\bar x\) be a \(k\)-tuple of
distinct variables, and let
\(\theta^+,\theta^-\in\mathrm{Pos}[\sigma\cup\{R\}]\) have free variables
among \(\bar x\) and parameters \(\bar p\). Assume that on every finite
\((\sigma\cup\{R\})\)-structure \(\theta^-\) is equivalent to
\(\neg\theta^+\), and that on every finite \(\sigma\)-structure, for all
values of \(\bar p\), the operator
\(F(X)\defeq\{\bar a:(\mathfrak A,X)\entails\theta^+(\bar a)\}\) is
monotone. The symbol \(R\) may occur in \(\theta^\pm\) in literals of both
signs. Then there are formulas
\(\mathrm{le}(\bar u,\bar w)\), \(\mathrm{lt}(\bar u,\bar w)\),
\(\mathrm{in}(\bar u)\) and \(\mathrm{out}(\bar u)\) in
\(\mathrm{Pos}[\sigma]\), with free variables among \(\bar u,\bar w,\bar p\)
for fresh \(k\)-tuples \(\bar u,\bar w\), which on every finite
\(\sigma\)-structure define \(\le_F\), \(<_F\),
\(\operatorname{lfp}(F)\) and \(A^k\setminus\operatorname{lfp}(F)\).
\end{hblemma}

The relations \(\le_F\) and \(<_F\) are defined by one simultaneous
induction whose clauses are a variant of Immerman's equations
\cite[p.~98]{Immerman1986}, written so that the stage
\(X^{|\bar c|_F}\) is named by \(\{\bar u:\bar u\le_F\bar c\}\) and its
complement by \(\{\bar u:\bar c<_F\bar u\}\). Both names are positive in
the relations being defined, which is why \(R\) may occur with both signs
in \(\theta^\pm\). The complement of the fixed point is expressed through a
tuple of maximal rank, Immerman's \(\mathrm{MAX}_\varphi\)
\cite[Lemmas~4.5--4.6, pp.~98--99]{Immerman1986}: a tuple lies outside
\(\operatorname{lfp}(F)\) exactly when the fixed point is empty or some
tuple of maximal rank has smaller rank than it.

\begin{hbtheorem}{thm:p03-w03-p03-059-immerman-normal-form}{proved=annex:p03-w03-p03-059-immerman-normal-form}{Immerman's normal form on finite structures}
\begin{enumerate}
 \item For every LFP formula \(\chi\) of a relational vocabulary
       \(\sigma\) there are formulas \(\chi^+,\chi^-\in\mathrm{Pos}[\sigma]\),
       with free variables among those of \(\chi\), such that on every
       finite \(\sigma\)-structure \(\chi^+\) is equivalent to \(\chi\) and
       \(\chi^-\) to \(\neg\chi\).
 \item Every formula of \(\mathrm{Pos}[\sigma]\) is equivalent on finite
       \(\sigma\)-structures to a formula of \(\mathrm{LFP}_1\) of
       vocabulary \(\sigma\).
 \item Consequently, on finite structures every LFP formula, and the
       negation of every LFP formula, is equivalent to a formula of
       \(\mathrm{LFP}_1\). In particular the complement of an
       \(\mathrm{LFP}_1\)-definable query on finite structures is again
       \(\mathrm{LFP}_1\)-definable, and by
       \cref{thm:p03-w03-p03-059-gfp-dual} every greatest fixed point is
       expressible without negated fixed points.
\end{enumerate}
Item~3 is Theorem~3.3.51 of \cite[p.~176]{Gradel2007FMTDescriptiveChapter}, stated there
without proof, with a pointer to stage comparison and to Immerman. It is
the content of Immerman's Theorem~2, that the fixed-point hierarchy
collapses at its first level \cite[pp.~87, 94]{Immerman1986}.
\end{hbtheorem}

The annex proof follows Immerman's sequence of lemmas. Item~2 combines
Immerman's simultaneous induction lemma (two simultaneous inductions are
one, with a pair of equal or unequal elements as a tag) and his
transitivity lemma (a fixed point inside the body of another fixed point
can be computed in a joint induction)
\cite[Lemmas~4.1--4.2 and Cor.~4.3, pp.~95--97]{Immerman1986}. Immerman
refers to Moschovakis for the inductions behind both; the annex writes them
out, the second in the form of Beki\'c's principle. Item~1 uses
\cref{thm:p03-w03-p03-059-stage-comparison} for least fixed points and, for
greatest fixed points, the dual operator. Immerman records a difficulty
raised by a referee: once negations are removed, the variable of an outer
fixed point may occur both positively and negatively inside an inner one
\cite[Lemma~4.7, p.~99]{Immerman1986}. He resolves it by substituting
names of the outer stage and of its complement, which are positive
\cite[p.~100]{Immerman1986}. Item~1 builds that remedy into the induction:
each formula is carried together with a positive form of its negation, and
\cref{thm:p03-w03-p03-059-stage-comparison} accepts a fixed-point variable
occurring with both signs. Two differences of setting remain. Immerman
assumes a constant symbol in every vocabulary and proves the sharper form
in which a single application of the fixed-point operator suffices
\cite[pp.~89, 94]{Immerman1986}; the statement above is Gr\"adel's
\(\mathrm{LFP}_1\) form, for relational vocabularies without constants. The
pair tag needs two distinct elements, so the annex treats one-element
structures separately, by unfolding the induction.

\begin{hbexample}{def:p03-w03-p03-059-non-reachability}{Non-reachability without a negated fixed point}
Let \(\theta^+(T,x,y)\defeq Exy\lor\exists z(Exz\land Tzy)\), the body of
\cref{def:p03-w03-p03-059-tc}, and
\(\theta^-(T,x,y)\defeq\neg Exy\land\forall z(\neg Exz\lor\neg Tzy)\). On a
finite graph the rank of a pair \((a,b)\) is the number of edges of a
shortest nonempty path from \(a\) to \(b\), and \(\infty\) if there is
none. A pair \((c,d)\) has maximal rank when no shortest path between any
pair is longer than the one from \(c\) to \(d\). The formula
\(\mathrm{out}(u,v)\) of \cref{thm:p03-w03-p03-059-stage-comparison} says:
either there are no edges at all, or some pair \((c,d)\) of maximal rank
has smaller rank than \((u,v)\). In the second case there is no path from
\(u\) to \(v\), because every pair joined by a path has rank at most that
of \((c,d)\). Maximality of \((c,d)\) is written
\[
 \mathrm{le}(c,d,c,d)\land\forall x\forall y\bigl(\mathrm{le}(x,y,c,d)\lor
 \theta^-(x,y)[T:\mathrm{le}(\cdot,\cdot,c,d)/\mathrm{lt}(c,d,\cdot,\cdot)]\bigr):
\]
every pair has rank at most that of \((c,d)\) or does not enter at the next
stage. For the path \(a\to b\to e\) with edges \(Eab\) and \(Ebe\), the
ranks are \(|(a,b)|=|(b,e)|=1\) and \(|(a,e)|=2\); the pair \((a,e)\) is the
only pair of maximal rank, and \((e,a)\) satisfies \(\mathrm{out}\) because
\(|(a,e)|=2<\infty=|(e,a)|\). No negation stands in front of a fixed point:
the negated \(T\)-atoms of \(\theta^-\) have been replaced by positive
occurrences of \(\mathrm{lt}\), and the only negated atoms left are
\(E\)-atoms.
\end{hbexample}

\begin{hbwarning}{def:p03-w03-p03-059-normal-form-boundary}{The normal form is a finite-structure theorem}
The hypothesis ``finite structures'' cannot be dropped. Gr\"adel reports
that on infinite structures some LFP formulas are equivalent to no positive
formula, and that the hierarchy obtained by alternating least and greatest
fixed points is strict there \cite[p.~176]{Gradel2007FMTDescriptiveChapter}; he refers to
work of Bradfield and of Moschovakis for those results. The normal form
therefore does not say that LFP is closed under complement as a matter of
logic: it is a consequence of finiteness. Immerman makes the same remark
\cite[p.~100]{Immerman1986}. The proof shows where finiteness enters. The
clauses of the stage comparison induction reach a tuple of rank \(r\) through
a tuple of rank \(r-1\), and on an infinite structure a stage may have no
immediate predecessor, which Immerman notes \cite[pp.~97--98]{Immerman1986};
and the negation of a fixed point is expressed through a tuple of maximal
rank, which exists because the closure ordinal is finite. The normal form
also concerns formulas over the stated vocabulary only; no order is needed
for it.
\end{hbwarning}

\section{Partial fixed-point logic}
\label{sec:p03-w03-p03-059-pfp}

Positivity guarantees that the stages of an operator increase. Without it,
the stages may oscillate. On a finite structure there are only finitely
many relations of a given arity, so an oscillating sequence must eventually
repeat. Partial fixed-point logic uses this to give every formula a
fixed-point reading on finite structures.

\begin{hbdefinition}{def:p03-w03-p03-059-pfp}{Partial fixed points and PFP}
Let \(\mathfrak A\) be a finite \(\tau\)-structure and \(\psi(R,\bar x)\) an
arbitrary formula of vocabulary \(\tau\cup\{R\}\), with \(R\) of arity
\(k\) and the remaining free variables fixed by an assignment. Its operator
\(F_\psi(R)\defeq\{\bar a:(\mathfrak A,R)\entails\psi(R,\bar a)\}\) need not
be monotone. Its \emph{finite stages} are
\[
 R^0\defeq\varnothing,\qquad R^{m+1}\defeq F_\psi(R^m)\quad(m\in\N).
\]
The sequence need not increase. If some stage satisfies
\(R^{m+1}=R^m\), the sequence has \emph{reached a fixed point}, and the
\emph{partial fixed point} of \(F_\psi\) is that stage \(R^m\) (it does not
depend on the choice of such \(m\), since from \(m\) on every stage equals
\(R^m\)). Otherwise the partial fixed point is the empty relation.

PFP extends first-order logic by the \emph{partial-fixed-point formation
rule}: from any formula \(\psi(R,\bar x)\), a \(k\)-tuple \(\bar x\) of
distinct variables and a \(k\)-tuple \(\bar t\) of terms, form
\[
 [\mathbf{pfp}\,R\bar x\,.\,\psi(R,\bar x)](\bar t),
\]
with free variables and parameters as in \cref{def:p03-w03-p03-059-lfp}. On
a finite structure it holds when \(\bar t^{\mathfrak A}\) belongs to the
partial fixed point of \(F_\psi\). The formula \(\psi\) may contain
negations of \(R\), quantifiers, and nested \(\mathbf{pfp}\) formulas. This
is the presentation of \cite[p.~177]{Gradel2007FMTDescriptiveChapter}.
\index{partial fixed point}\index{partial fixed-point logic}\index{PFP}
\end{hbdefinition}

\begin{hblemma}{thm:p03-w03-p03-059-pfp-dichotomy}{proved=annex:p03-w03-p03-059-pfp-dichotomy}{The two possible behaviours of finite stages}
Let \(\mathfrak A\) be finite with \(n\) elements, \(F_\psi\) as in
\cref{def:p03-w03-p03-059-pfp}, and \(N\defeq2^{n^k}\). Exactly one of
the following holds:
\begin{enumerate}
 \item the stages reach a fixed point, and then they do so at some
       \(m<N\);
 \item the stages are eventually periodic with least period greater than
       one: there are \(i<j\leq N\) with \(R^i=R^j\), and no stage is a
       fixed point.
\end{enumerate}
Consequently the partial fixed point equals \(R^N\) when
\(F_\psi(R^N)=R^N\), and equals \(\varnothing\) otherwise
\cite[p.~177]{Gradel2007FMTDescriptiveChapter}.
\end{hblemma}

\begin{hbtheorem}{thm:p03-w03-p03-059-lfp-in-pfp}{proved=annex:p03-w03-p03-059-lfp-in-pfp}{LFP is contained in PFP on finite structures}
If \(R\) occurs only positively in \(\psi(R,\bar x)\), then on every finite
structure
\[
 [\mathbf{lfp}\,R\bar x\,.\,\psi(R,\bar x)](\bar t)
 \;\Leftrightarrow\;
 [\mathbf{pfp}\,R\bar x\,.\,\psi(R,\bar x)](\bar t).
\]
Hence every LFP formula is equivalent on finite structures to a PFP
formula; in Gr\"adel's notation \(\mathrm{LFP}\leq\mathrm{PFP}\)
\cite[p.~177]{Gradel2007FMTDescriptiveChapter}.
\end{hbtheorem}

The inclusion has a quantitative side. By
\cref{thm:p03-w03-p03-059-finite-stages}, a least-fixed-point iteration on
a structure with \(n\) elements stops after at most \(n^k\) steps. A
partial-fixed-point iteration can take exponentially many.

\begin{hbexample}{def:p03-w03-p03-059-counter}{A binary counter as a partial-fixed-point induction}
On a finite linear order \((A,<)\) with \(n\geq1\) elements, Gr\"adel's
Example~3.3.53 uses the formula
\[
 \psi(R,x)\defeq
 \bigl(Rx\land\exists y(y<x\land\neg Ry)\bigr)\lor
 \bigl(\neg Rx\land\forall y(y<x\to Ry)\bigr)\lor
 \forall y\,Ry
\]
\cite[p.~178]{Gradel2007FMTDescriptiveChapter}. Read a set \(S\subseteq A\) as a binary
numeral: the element of rank \(j\) in the order (rank \(0\) for the least
element) carries the bit of weight \(2^j\), and the bit is \(1\) when the
element lies in \(S\). The first disjunct keeps a \(1\) bit when some lower
bit is \(0\); the second turns a \(0\) bit into \(1\) when all lower bits
are \(1\); together they add one to the numeral. The third disjunct makes
the full set \(A\), the numeral \(2^n-1\), a fixed point instead of letting
it overflow to \(\varnothing\). So \(R^m\) is the numeral \(m\), and the
iteration visits all \(2^n\) subsets before it stops.
\cref{thm:p03-w03-p03-059-counter} states and proves this.

For \(n=3\) with \(a<b<c\), the stages \(R^0,\ldots,R^7\) are
\[
 \varnothing,\ \{a\},\ \{b\},\ \{a,b\},\ \{c\},\ \{a,c\},\ \{b,c\},\ \{a,b,c\},
\]
and \(\{a,b,c\}\) is the first stage with \(F_\psi(R^m)=R^m\), at \(m=7\). The
\(R\) occurring negatively in the first and second disjuncts makes
\(F_\psi\) non-monotone, so this is not an LFP induction: for instance
\(\{a\}\subseteq\{a,b\}\), but \(F_\psi(\{a\})=\{b\}\) is not contained
in \(F_\psi(\{a,b\})=\{c\}\).
\end{hbexample}

\begin{hblemma}{thm:p03-w03-p03-059-counter}{proved=annex:p03-w03-p03-059-counter}{The counter reaches its fixed point at stage \(2^n-1\)}
For \(\psi\) and \((A,<)\) as in \cref{def:p03-w03-p03-059-counter}, with
\(n=|A|\geq1\), the stages are pairwise distinct for \(m\leq2^n-1\),
\(R^{2^n-1}=A\), and \(F_\psi(A)=A\). Hence the partial fixed point is
\(A\), first reached at stage \(2^n-1\)
\cite[p.~178]{Gradel2007FMTDescriptiveChapter}.
\end{hblemma}

\subsection{Complexity of partial fixed points}
\label{sec:p03-w03-p03-059-pfp-complexity}

Exponentially many stages need not cost exponential space: only the current
stage and a counter must be stored. To speak of complexity at all,
Gr\"adel first fixes how an ordered structure is presented to a machine.

\begin{hbdefinition}{def:p03-w03-p03-059-good-encoding}{Good encodings of ordered structures}
For a finite relational vocabulary \(\tau\), let \(\mathrm{Ord}(\tau)\) be
the class of structures \((\mathfrak A,<)\) with \(\mathfrak A\) a finite
\(\tau\)-structure and \(<\) a linear order of its universe. For
\((\mathfrak A,<)\) with \(n\) elements, \(A^k\) is identified with
\(\{0,\ldots,n^k-1\}\) by the rank of each tuple in the lexicographic order
induced by \(<\). An encoding \(\mathit{code}\) of \(\mathrm{Ord}(\tau)\) by
words over a finite alphabet \(\Sigma_0\) is \emph{good} when:
\begin{enumerate}
 \item[(i)] \(\mathit{code}(\mathfrak A,<)=\mathit{code}(\mathfrak B,<)\)
       if and only if the two ordered structures are isomorphic;
 \item[(ii)] \(|\mathit{code}(\mathfrak A,<)|\leq p(|A|)\) for a fixed
       polynomial \(p\);
 \item[(iii)] for every \(k\) and every \(\sigma\in\Sigma_0\) there is a
       first-order formula \(\beta_\sigma(x_1,\ldots,x_k)\) of vocabulary
       \(\tau\cup\{<\}\) such that \((\mathfrak A,<)\entails
       \beta_\sigma(\bar a)\) exactly when the symbol at position
       \(\bar a\) of \(\mathit{code}(\mathfrak A,<)\) is \(\sigma\);
 \item[(iv)] from \(\mathit{code}(\mathfrak A,<)\), a relation symbol and a
       tuple, the truth of the atomic statement can be decided
       efficiently.
\end{enumerate}
This is Definition~3.1.10 of \cite[p.~136]{Gradel2007FMTDescriptiveChapter}. In (iii),
positions beyond the end of the word carry no symbol, so every
\(\beta_\sigma\) is false there. A class of ordered structures is
\emph{recognizable} in a complexity class when the set of its codes is.

The complexity statements below also use three properties that the source
takes for granted when it passes between structures and their codes, and
which are stated here as hypotheses: (a) the set of all codes is decidable
in polynomial time, and from a code one computes \(|A|\) in polynomial
time, the atomic test of (iv) also running in polynomial time; (b) from
\(|A|\) and the tables of the relations of an ordered structure, its code
is computable in polynomial time; (c)
\(|\mathit{code}(\mathfrak A,<)|\geq|A|\). Gr\"adel's convenient encoding
\cite[p.~136]{Gradel2007FMTDescriptiveChapter} begins with \(|A|\) ones, followed by
padding zeros and the characteristic strings of the relations in the
lexicographic tuple order; its length is a fixed polynomial in \(|A|\)
determined by the leading run of ones, which gives (a)--(c).
\index{good encoding}
\end{hbdefinition}

\begin{hbtheorem}{thm:p03-w03-p03-059-pfp-pspace}{proved=annex:p03-w03-p03-059-pfp-pspace}{PFP formulas are evaluated in polynomial space}
For every PFP sentence \(\psi\), the class of finite models of \(\psi\) is
in \(\mathsf{PSPACE}\): for an encoding as in
\cref{def:p03-w03-p03-059-good-encoding}, satisfying (a)--(c) there, a
deterministic algorithm decides, given a word, whether it is a code
\(\mathit{code}(\mathfrak A,<)\) of some finite structure with
\(\mathfrak A\entails\psi\), using work space polynomial in the length of
the word. The same holds for a PFP formula \(\psi(\bar x)\) with free
first-order variables, whose finite models are the pairs
\((\mathfrak A,\bar a)\): given a code and the ranks of \(\bar a\) written in
binary, the truth of \(\psi(\bar a)\) is decided in work space polynomial in
the length of the input. In short, \(\mathrm{PFP}\subseteq\mathsf{PSPACE}\).
This is Proposition~3.3.54 of \cite[p.~178]{Gradel2007FMTDescriptiveChapter}, which is
stated there for every PFP formula.
\end{hbtheorem}

\begin{hbtheorem}{thm:p03-w03-p03-059-pfp-captures}{proved=annex:p03-w03-p03-059-pfp-captures}{Abiteboul, Vianu and Vardi: PFP captures PSPACE on ordered structures}
Fix a finite relational vocabulary \(\tau\) and a good encoding of
\(\mathrm{Ord}(\tau)\) with properties (a)--(c) of
\cref{def:p03-w03-p03-059-good-encoding}. A class \(\mathcal K\subseteq\mathrm{Ord}(\tau)\),
closed under isomorphism, is recognizable in \(\mathsf{PSPACE}\) if and only
if it is the class of finite ordered models of a PFP sentence of vocabulary
\(\tau\cup\{<\}\). This is Theorem~3.3.55 of
\cite[p.~178]{Gradel2007FMTDescriptiveChapter}, where the order is part of every input
structure.
\end{hbtheorem}

\begin{hbcorollary}{thm:p03-w03-p03-059-ptime-pspace}{proved=annex:p03-w03-p03-059-ptime-pspace}{\(\mathsf P=\mathsf{PSPACE}\) as a question about fixed points}
The following are equivalent:
\begin{enumerate}
 \item \(\mathsf P=\mathsf{PSPACE}\);
 \item on ordered finite structures LFP and PFP have the same expressive
       power: for every finite relational \(\tau\), a class of ordered
       finite \(\tau\)-structures is defined by an LFP sentence of
       vocabulary \(\tau\cup\{<\}\) if and only if it is defined by a PFP
       sentence of that vocabulary.
\end{enumerate}
This is the consequence stated at the head of \cite[p.~180]{Gradel2007FMTDescriptiveChapter}.
The proof uses \cref{thm:p03-w03-p03-059-pfp-captures} together with the
capture of \(\mathsf P\) by LFP on ordered finite structures
\cite[Theorem~3.3.16, p.~157]{Gradel2007FMTDescriptiveChapter}, which the handbook states
as \cref{thm:p04-immerman-vardi} and proves in Immerman's syntax as
\cref{thm:p09-immerman1986-capture}.
\end{hbcorollary}

The corollary concerns ordered structures. For several pairs of logics a
separation on arbitrary finite structures is easy even though the
separation on ordered structures would settle an open complexity question;
Gr\"adel's examples are deterministic against general transitive-closure
logic and transitive-closure logic against LFP
\cite[p.~180]{Gradel2007FMTDescriptiveChapter}. For LFP against PFP the unordered question
is no easier than the ordered one.

\begin{hbtheorem}{thm:p03-w03-p03-059-abiteboul-vianu}{proved=annex:p03.av.4}{Abiteboul--Vianu equivalence}
LFP and PFP have the same expressive power on finite structures if
and only if \(\mathsf P=\mathsf{PSPACE}\). More precisely, under
that complexity-class equality, every PFP formula over a finite
vocabulary has an equivalent LFP formula over the same vocabulary,
with free first-order variables among those of the original formula.
Equivalence holds on all finite structures, without a distinguished
order on their universes.
\end{hbtheorem}

Gr\"adel states this as Theorem~3.3.56 without a proof
\cite[p.~180]{Gradel2007FMTDescriptiveChapter}.
The construction in \cref{sec:p03-abiteboul-vianu-completion}
proves the unordered implication by ordering tuple types and then
pulling the resulting LFP definition back to the original vocabulary.
The reverse implication uses the ordered capture results of
\cref{thm:p03-w03-p03-059-ptime-pspace}.

\begin{hbwarning}{def:p03-w03-p03-059-pfp-limits}{What the partial-fixed-point results do not say}
\begin{enumerate}
 \item \emph{Finite structures only.} The partial fixed point is defined
       through the finite stages \(R^m\), \(m\in\N\), on a finite
       structure (\cref{def:p03-w03-p03-059-pfp}). On an infinite
       structure the stage sequence need not repeat and the definition
       assigns no value, so neither \cref{thm:p03-w03-p03-059-lfp-in-pfp}
       nor \cref{thm:p03-w03-p03-059-abiteboul-vianu} says anything about
       infinite structures. Both LFP and PFP are compared there only as
       logics of finite structures.
 \item \emph{Cycling and non-membership look alike.} The value
       \(\varnothing\) is returned both when the stages cycle and when the
       iteration stops at the empty relation; a formula
       \(\neg[\mathbf{pfp}\,R\bar x\,.\,\psi](\bar t)\) cannot tell which
       happened.
 \item \emph{Order is a hypothesis of the capture theorems.}
       \cref{thm:p03-w03-p03-059-pfp-captures,thm:p03-w03-p03-059-ptime-pspace}
       are about ordered structures with the order in the vocabulary. The
       unordered equivalence of \cref{thm:p03-w03-p03-059-abiteboul-vianu}
       is a separate theorem, and the separation of LFP from PFP on all
       finite structures is exactly as hard as separating
       \(\mathsf P\) from \(\mathsf{PSPACE}\).
 \item \emph{Simultaneous inductions.} Gr\"adel notes that simultaneous
       partial-fixed-point inductions reduce to simple ones only at the cost
       of raising the arity, and that it is unknown whether simultaneous
       monadic PFP collapses to simple monadic PFP
       \cite[p.~178]{Gradel2007FMTDescriptiveChapter}. The proof of
       \cref{thm:p03-w03-p03-059-pfp-captures} given here uses one
       relation of higher arity in place of a simultaneous induction.
\end{enumerate}
\end{hbwarning}

\section{Relationships and sources}
\label{sec:p03-w03-p03-059-relationships}

The first-order least fixed-point constructor of \cref{def:p04-fol-lfp} and
the ordered capture of \(\mathsf P\) (\cref{thm:p04-immerman-vardi}, proved
in Immerman's syntax as \cref{thm:p09-immerman1986-capture}) are the LFP
side of \cref{thm:p03-w03-p03-059-ptime-pspace}. The same lattice-theoretic
argument as \cref{thm:p03-w03-p03-059-knaster-tarski} gives the least and
greatest fixed points of the modal \(\mu\)-calculus
(\cref{thm:p08-bvb501-mu-constructions}), where the lattice is the power set
of the set of worlds. The zero--one law for fixed-point logics recorded in
\cref{sec:p09-presentation-n-c-002} covers PFP because PFP is contained in
the finite-variable infinitary logic used there; the order needed in
\cref{thm:p03-w03-p03-059-pfp-captures} is exactly what destroys that law.

\paragraph{Exact sources.} Operator theory: Definition~3.3.1,
Theorem~3.3.2, Lemma~3.3.3, Proposition~3.3.4, Exercise~3.3.5 and
Lemma~3.3.6 of \cite[pp.~153--155]{Gradel2007FMTDescriptiveChapter}. LFP syntax and the
transitive-closure example: Definition~3.3.7 and Example~3.3.8 of
\cite[p.~155]{Gradel2007FMTDescriptiveChapter}; greatest fixed points through negation:
\cite[p.~156]{Gradel2007FMTDescriptiveChapter}; \(\mathrm{LFP}_1\):
\cite[p.~173]{Gradel2007FMTDescriptiveChapter}; Immerman's normal form and its
infinite-structure boundary: Theorem~3.3.51 of
\cite[p.~176]{Gradel2007FMTDescriptiveChapter}; simultaneous fixed points:
Theorem~3.3.43 of \cite[p.~174]{Gradel2007FMTDescriptiveChapter}. The proof of the normal
form: Immerman's Theorem~2 \cite[pp.~87, 94]{Immerman1986}, his constant
assumption \cite[p.~89]{Immerman1986}, Lemmas~4.1 and~4.2 with
Corollary~4.3 \cite[pp.~95--97]{Immerman1986}, rank, stage comparison and
Fact~4.4 \cite[p.~97]{Immerman1986}, the stage comparison equations and
Lemma~4.5 \cite[p.~98]{Immerman1986}, Lemmas~4.6 and~4.7
\cite[pp.~98--99]{Immerman1986}, and the closing argument and the remark
that the result fails on infinite structures
\cite[p.~100]{Immerman1986}. PFP, the inclusion of LFP, and the
counter: \cite[pp.~177--178]{Gradel2007FMTDescriptiveChapter}, Example~3.3.53;
Proposition~3.3.54 and Theorem~3.3.55 with its proof:
\cite[pp.~178--179]{Gradel2007FMTDescriptiveChapter}; good encodings:
Definition~3.1.10 of \cite[p.~136]{Gradel2007FMTDescriptiveChapter}; the
\(\mathsf P\) versus \(\mathsf{PSPACE}\) consequence and
Theorem~3.3.56: \cite[p.~180]{Gradel2007FMTDescriptiveChapter}.
